\chapter{Discourse}\label{ch:discourse}

% ============================================================
% STATUS: SCAFFOLD — largely new; no direct predecessor in 2020 archive.
% Some material in archive/2020/chapters/08-semantics-and-usage.tex may be relevant.
% This chapter is theoretically the most ambitious of the rewrite:
% it reframes the alignment system as a discourse-level phenomenon and
% treats the converb chain as the basic unit of narrative structure.
% ============================================================

\section{Sentence-chains as discourse units}\label{sec:discourse-chains}

% ============================================================
% Iridian discourse is organized in "sentence chains": sequences of
% converb clauses capped by a finite verb. Each sentence chain functions
% as a single discourse unit. The chain has:
%   - Converb clauses = background, supporting, or subordinate events
%   - Finite verb = the foregrounded, topically relevant event
% This parallels the foreground/background distinction found in many
% verb-final or SOV languages with converb systems.
% ============================================================

\section{Topic continuity and topic shift}\label{sec:topic-continuity}

% The absolutive participant tends to persist as topic across sentence
% chains. When the topic shifts, this is typically signaled by:
%   - Lexical re-introduction of the new topic
%   - Switch-reference clitic on the converb
%   - A demonstrative or proximate/obviative distinction
% The obviative/proximate system of demonstratives (ša vs. dní) is the
% primary cross-clause topic-tracking device.

\section{Definiteness in discourse context}\label{sec:discourse-definiteness}

% The definiteness marker -ot/-t on a noun signals that the referent is
% identifiable to the hearer. In discourse, first mention is typically
% bare (no definiteness marker); subsequent mention may use the marker,
% the demonstrative, or drop the noun entirely (pro-drop).
% The interaction of definiteness, topic, and absolutive alignment
% is the key thread here.

\section{Aspect choice as narrative strategy}\label{sec:aspect-narrative}

% The split-ergative alignment means that aspect choice simultaneously
% encodes:
%   (a) the temporal/aspectual shape of the event
%   (b) which participant is the discourse topic
% In narrative, speakers use the perfective/retrospective to bring the
% patient into topic/absolutive position (foregrounding the event outcome)
% and use nominative aspects to keep the agent as topic (foregrounding
% the agent's activity).
% This is the core pragmatic principle behind the aspect system.

In extended discourse, aspect choice is inseparable from participant tracking. The perfective and retrospective place the patient in the absolutive and are therefore preferred when the discourse is about what happened \emph{to} a participant: they are the natural vehicle for narrative foreground, event culmination, and topic shift to an affected referent. The nominative aspects, by contrast, preserve the agent or sole participant as absolutive and are accordingly favoured where the continuity of an actor matters more than the bounded result of the event.

The same principle explains the discourse value of the antipassive. Speakers choose the antipassive when a transitive event is relevant chiefly as part of an agent's ongoing activity and not because the patient itself needs to be established, identified, or tracked. The construction keeps the agent in absolutive position, suppresses any patient from the grammatical core, and thereby functions as a topic-maintenance device. It is especially natural when the patient is generic, incorporated into the event type, inferable from context, or simply of no further narrative consequence; for the morphological details see Chapter~\ref{ch:verbs}.

The causative serves a different discourse purpose. By adding a causer, it allows the speaker to recast an event as one of inducement, delegated action, permission, or coercion. This is a useful viewpoint-management device in narrative: a simple predicate presents an event as directly performed, whereas a causative can foreground authority, distribute responsibility across two participants, or distance the discourse topic from the event's direct execution. In perfective contexts this often yields a strongly interventionist reading; in retrospective contexts it more readily supports permissive, indirect, or weak-control interpretations. The formal argument structure of the causative is described in Chapter~\ref{ch:verbs}; what matters here is that the construction lets the speaker decide whether the discourse should centre on the doing, the causing, or the managed consequences of the action.

\section{Tail-head linkage}\label{sec:tail-head}

% Tail-head linkage is a discourse-structuring device where the final
% clause of one sentence chain is repeated (in whole or in abbreviated
% form) at the start of the next chain. This is attested in the sample
% texts and is a common device in narrative Iridian.

\section{Discourse particles}\label{sec:discourse-particles}

% ============================================================
% STATUS: STUB — to be documented
% Categories:
%   - Topic-management particles (topic shift, topic continuation)
%   - Stance particles (mirative, dubitative, concessive)
%   - Connectives (now, then, well, anyway)
% Some of these were described as "evidential" in the 2020 archive
% (under the QUOT/Infer/Rep labels). They should be reconciled with
% the evidential slot in the morphological template.
% ============================================================

\section{Register variation}\label{sec:register-discourse}

% Literary vs. colloquial clause-linking preferences:
%   - Literary: tends toward full converb chains, explicit nominalization
%   - Colloquial: nu as discourse particle, purposive-future, shorter chains
%   - The increasing use of finite coordinated clauses under Slavic influence
